%% Example CV - Rania Boucenna
%% Comprehensive overview of all competencies, experience, and achievements.
%% Use as a master reference when tailoring role-specific CVs.
%%
%% Compile with: lualatex main_example.tex
%% (pdflatex often fails on modern MiKTeX with fontawesome5 font-expansion errors.)

\documentclass[11pt,a4paper,sans]{moderncv}
\moderncvstyle{banking}
\moderncvcolor{blue}

% Force both first and last name AND section headings to render in moderncv
% blue (color1). Default banking on lualatex+MiKTeX leaves these black, which
% looks inconsistent with the rest of the blue accent scheme.
\renewcommand*{\firstnamestyle}[1]{{\fontsize{34}{36}\bfseries\upshape\color{color1}#1}}
\renewcommand*{\lastnamestyle}[1]{{\fontsize{34}{36}\bfseries\upshape\color{color1}#1}}
\renewcommand*{\sectionstyle}[1]{{\sectionfont\color{color1}#1}}

\usepackage[utf8]{inputenc}
\usepackage{hyperref}
\hypersetup{
    colorlinks=true,
    linkcolor=blue,
    filecolor=magenta,
    urlcolor=blue,
    pdftitle={Rania Boucenna - CV},
    pdfpagemode=FullScreen,
}
\usepackage[scale=0.80]{geometry}
\usepackage{import}

% personal data
\name{Rania}{Boucenna}
\address{Rouen, France}{}{}
\phone[mobile]{+33 7 59 08 89 32}
\email{raniaboucenna3@gmail.com}
\extrainfo{\href{https://www.linkedin.com/in/raniaboucenna}{LinkedIn}}

\begin{document}

\makecvtitle

% ============================================================
%     PROFILE STATEMENT
% ============================================================

\vspace{6pt}
\small{Ingénieure en génie des procédés (École Nationale Polytechnique de Constantine), actuellement en Master 2 Procédés Innovants et Développement Durable à l'Université Sorbonne Paris-Nord. Expérience en R\&D sur les matériaux biosourcés, la caractérisation physico-chimique et l'analyse de données sous Python, ainsi qu'en génie des procédés industriels (raffinage, traitement des eaux, pharma). Méthodique et autonome, je vise à optimiser l'efficacité des procédés et à concevoir des solutions innovantes et écoresponsables.}

% ============================================================
%     CORE COMPETENCIES
% ============================================================

\section{Compétences Clés}
\vspace{1pt}
\begin{itemize}

\item \textbf{Génie des procédés}: Bilans matière et énergie, dimensionnement et optimisation des procédés, chimie des réacteurs

\item \textbf{Valorisation de la biomasse \& matériaux biosourcés}: Élaboration et caractérisation de matériaux biosourcés (béton biosourcé, bio-patchs)

\item \textbf{Caractérisation physico-chimique}: GCMS, UV-visible, FTIR, granulométrie laser

\item \textbf{Outils de simulation \& données}: Aspen HYSYS, Python, Matlab, Fluent

\item \textbf{Développement durable \& QHSE}: Économie circulaire, démarche QHSE

\end{itemize}

% ============================================================
%     PROFESSIONAL EXPERIENCE
% ============================================================

\section{Expérience Professionnelle}
\vspace{3pt}
\begin{itemize}

% --- Most Recent Role ---
\item{\cventry{2026-Présent}{Ingénieure en développement durable (stage de fin d'études)}{Laboratoire Eclore, UniLaSalle Rouen}{Rouen, France}{}{\vspace{1pt}
\begin{itemize}
    \item Mise au point d'un béton biosourcé destiné à l'isolation thermique
    \item Fabrication de matériaux isolants biosourcés
    \item Caractérisation physico-chimique, mécanique et hygroscopique des matériaux élaborés
\end{itemize}}}

\vspace{3pt}

% --- Previous Role ---
\item{\cventry{Jan 2023 - Jul 2023}{Ingénieure R\&D - Projet de Fin d'Études}{École Nationale Polytechnique de Constantine}{Constantine, Algérie}{}{\vspace{1pt}
\begin{itemize}
    \item Recherche sur la gélation ionotropique des huiles essentielles pour concevoir des bio-patchs de cicatrisation
    \item Encapsulation ionotropique par pompe à seringue programmable ; caractérisation (GCMS, UV-visible, FTIR, granulométrie laser)
    \item Modélisation mathématique de la cinétique de libération et optimisation des paramètres de formulation
    \item Mention Très Bien
\end{itemize}}}

\vspace{3pt}

% --- Earlier Role ---
\item{\cventry{Jul 2022 - Aug 2022}{Ingénieure process}{Sonatrach, Raffinerie de Skikda}{Skikda, Algérie}{}{\vspace{1pt}
\begin{itemize}
    \item Acquisition du savoir-faire sur les étapes du raffinage du pétrole brut (dessalage, distillation, traitement des produits finis)
\end{itemize}}}

\end{itemize}

% ============================================================
%     EDUCATION
% ============================================================

\section{Formation}
\vspace{1pt}
\begin{itemize}

\item{\cventry{2025-2026}{Master 2, Procédés Innovants et Développement Durable}{Université Sorbonne Paris-Nord (Institut Galilée)}{Île-de-France, France}{}{\vspace{1pt}
En cours, soutenance prévue septembre 2026. Simulation des procédés, développement durable, QHSE.
}}

\vspace{3pt}

\item{\cventry{2020-2023}{Diplôme d'Ingénieure, Génie des Procédés}{École Nationale Polytechnique de Constantine (ENPC)}{Constantine, Algérie}{}{\vspace{1pt}
Thèse (PFE): ``Gélation ionotropique des huiles essentielles pour la conception de bio-patchs de cicatrisation des plaies.'' Mention Très Bien.
}}

\end{itemize}

% ============================================================
%     LANGUAGES
% ============================================================

\section{Langues}
\vspace{1pt}
\begin{itemize}
\item Français (native), Arabe (native), Anglais (C1), Allemand (débutant).
\end{itemize}

% ============================================================
%     PUBLICATIONS (optional)
% ============================================================

% None yet - section omitted from the master CV until a publication exists.

% ============================================================
%     HONORS AND AWARDS (optional)
% ============================================================

\section{Distinctions}
\vspace{1pt}
\begin{itemize}
\item{\textbf{Mention Très Bien} -- Projet de Fin d'Études, École Nationale Polytechnique de Constantine (2023).}
\end{itemize}

% ============================================================
%     REFERENCES
% ============================================================

\section{Références}
\vspace{1pt}
\begin{itemize}
\item{Disponibles sur demande.}
\end{itemize}

\end{document}
